\documentclass[11pt]{article}
\usepackage[a4paper,margin=1in]{geometry}
\usepackage{amsmath}
\usepackage{amssymb}
\usepackage{array}
\usepackage{booktabs}
\usepackage{hyperref}

\title{INP to Structural Graph: Exact Algorithm Walkthrough}
\author{}
\date{\today}

\begin{document}
\maketitle

\section{Scope}
This note explains exactly how the current code transforms an Abaqus \texttt{.inp} file into a structural graph.
It is implementation-aligned with:
\begin{itemize}
  \item \texttt{Matlab/GNN/datasetCreation/step3\_batch\_structural\_graph\_from\_inp.m}
  \item \texttt{Matlab/gnn\_prep\_spinodal/src/io/read\_abaqus\_inp.m}
  \item \texttt{Matlab/gnn\_prep\_spinodal/src/mesh/select\_region\_elements.m}
  \item \texttt{Matlab/gnn\_prep\_spinodal/src/graph/build\_full\_reference\_graph.m}
  \item \texttt{Matlab/gnn\_prep\_spinodal/src/graph/extract\_structural\_graph.m}
\end{itemize}

No method comparison is included here. This is only a ``how it works'' specification.

\section{Simple Explanation}
\subsection*{Pipeline in plain language}
\begin{enumerate}
  \item Read \texttt{sheet\_shell.inp} and parse nodes, shell elements, and element sets.
  \item Pick the target element region (usually the spinodal shell set).
  \item Build a dense mesh graph:
  \begin{itemize}
    \item graph nodes = selected Abaqus nodes,
    \item graph edges = unique shell edges from element connectivity,
    \item boundary nodes = nodes on edges used by exactly one element.
  \end{itemize}
  \item Rasterize dense XY coordinates to a native grid mask.
  \item Skeletonize each connected component with Zhang-Suen thinning.
  \item Convert skeleton pixels back to a graph (skeleton graph).
  \item Compress degree-2 chains into longer structural edges, controlled by \texttt{DetailLevel}.
  \item Run a local cleanup pass:
  \begin{itemize}
    \item contract very short edges,
    \item merge nearby high-degree junction pairs.
  \end{itemize}
  \item Attach thickness summaries to each structural edge using the full-node path.
  \item Save one \texttt{sample.mat} per run with the final structural graph.
\end{enumerate}

\subsection*{Simple meaning of the main graphs}
\begin{itemize}
  \item Dense graph $G_d$: directly from selected shell mesh topology.
  \item Skeleton graph $G_s$: one-node-thick centerline graph.
  \item Structural graph $G_r$: reduced graph after chain compression and cleanup.
\end{itemize}

\section{Detailed Explanation}

\subsection*{Notation}
\[
G_d = (V_d, E_d), \quad
G_s = (V_s, E_s), \quad
G_r = (V_r, E_r).
\]
Coordinates use physical XY taken from Abaqus node coordinates.

\subsection*{Default parameters used by Step 3}
In \texttt{step3\_batch\_structural\_graph\_from\_inp.m}, each run is processed with:
\begin{itemize}
  \item \texttt{ElsetName = 'SPINODAL\_SHELL'}
  \item \texttt{AutoDetectElset = true}
  \item \texttt{PatternPriority = \{'spinodal','top'\}}
  \item \texttt{DetailLevel = 0.30}
  \item \texttt{MinIslandNodes = 1}
\end{itemize}

\subsection*{1. Batch orchestration (\texttt{step3\_batch\_structural\_graph\_from\_inp.m})}
\paragraph{Simple}
The script finds run folders, skips already-processed runs, then parallelizes per-run graph extraction.

\paragraph{Detailed}
\begin{enumerate}
  \item Recursively discover run folders under \texttt{Matlab/GNN/data/raw/samples}.
  \item Keep folders that contain both \texttt{sheet.mat} and \texttt{mesh\_manifest.json}.
  \item For each run:
  \begin{itemize}
    \item expect \texttt{sheet\_shell.inp},
    \item skip if missing,
    \item skip if output \texttt{sample.mat} already exists.
  \end{itemize}
  \item Start a parallel pool and process pending runs in \texttt{parfor}.
  \item For each run: call parser, full-graph builder, structural extractor, then save \texttt{sample.mat}.
\end{enumerate}

\subsection*{2. INP parsing (\texttt{read\_abaqus\_inp.m})}
\paragraph{Simple}
The parser reads the text file block by block and extracts:
nodes, elements, and elsets.

\paragraph{Detailed}
\begin{enumerate}
  \item Read file line-by-line.
  \item Ignore empty lines and Abaqus comment lines starting with \texttt{**}.
  \item Detect keywords (\texttt{*NODE}, \texttt{*ELEMENT}, \texttt{*ELSET}) and parse attributes.
  \item In \texttt{NODE} mode:
  \begin{itemize}
    \item parse CSV values,
    \item first value is node label,
    \item next values are coordinates (up to XYZ; missing values padded with 0).
  \end{itemize}
  \item In \texttt{ELEMENT} mode:
  \begin{itemize}
    \item first value is element label,
    \item remaining values are node labels in element connectivity,
    \item store element type from keyword attribute \texttt{TYPE=...}.
  \end{itemize}
  \item Handle long Abaqus connectivity continuation:
  if a line starts with comma, append values to previous element connectivity.
  \item In \texttt{ELSET} mode:
  \begin{itemize}
    \item parse element labels into the active elset,
    \item if \texttt{GENERATE} is present, expand triples \((a,b,s)\) into sequence \(a:s:b\).
  \end{itemize}
  \item Deduplicate each elset label list with stable order.
\end{enumerate}

Output structure contains:
\begin{itemize}
  \item \texttt{node\_labels}, \texttt{node\_coords},
  \item \texttt{element\_labels}, \texttt{element\_connectivity}, \texttt{element\_types},
  \item \texttt{elset\_names}, \texttt{elset\_values}.
\end{itemize}

\subsection*{3. Region selection (\texttt{select\_region\_elements.m})}
\paragraph{Simple}
Pick which elements are ``in region'' before graph building.

\paragraph{Detailed}
\begin{enumerate}
  \item If explicit \texttt{ElsetName} is given, use it exactly (case-insensitive).
  \item Else if auto-detect is enabled, scan elset names in pattern order
  (default: \texttt{spinodal} then \texttt{top}), and choose first match.
  \item If no elset is selected, fall back to all elements.
  \item Keep only selected elements that exist in \texttt{inpData.element\_labels}.
  \item Build \texttt{region.node\_labels} as stable unique concatenation of selected element connectivity.
\end{enumerate}

\subsection*{4. Dense graph construction (\texttt{build\_full\_reference\_graph.m})}
\paragraph{Simple}
Convert selected shell elements into an undirected graph and mark boundary nodes.

\paragraph{Detailed}
\begin{enumerate}
  \item Map selected Abaqus node labels to coordinates in \texttt{inpData.node\_coords}.
  \item Convert each selected element connectivity from Abaqus labels to local node indices.
  \item Extract shell edges per element:
  for ordered element nodes \((n_1,\dots,n_k)\), create edges
  \((n_1,n_2),(n_2,n_3),\dots,(n_k,n_1)\).
  \item Sort each edge endpoint pair so edges are undirected.
  \item Remove self edges \((i,i)\).
  \item Concatenate all element edges into \texttt{all\_edges}.
  \item Compute \texttt{unique\_edges} and \texttt{edge\_use\_count}.
  \item Boundary rule:
  edge is boundary iff its use count is exactly 1.
  \item Boundary nodes are all nodes incident to at least one boundary edge.
\end{enumerate}

The dense graph contains:
\begin{itemize}
  \item \texttt{node\_coords}, \texttt{node\_labels},
  \item \texttt{edges\_local}, \texttt{graph} (MATLAB graph object),
  \item \texttt{boundary\_mask}, \texttt{boundary\_labels},
  \item selected region element metadata.
\end{itemize}

\subsection*{5. Rasterization to native XY grid (\texttt{extract\_structural\_graph.m})}
\paragraph{Simple}
Map dense nodes onto a 2D occupancy grid so morphology and thinning can run.

\paragraph{Detailed}
\begin{enumerate}
  \item Use only XY from \texttt{fullGraph.node\_coords(:,1:2)}.
  \item Cluster X values and Y values independently with \texttt{cluster\_axis\_values}:
  \begin{itemize}
    \item sort values,
    \item split groups when adjacent gap exceeds tolerance,
    \item tolerance is based on smallest positive spacing and machine-level fallback,
    \item use group means as grid axis values.
  \end{itemize}
  \item For each full node:
  \begin{itemize}
    \item get bin \((r,c)\),
    \item write first node index into \texttt{nodeAt(r,c)}.
  \end{itemize}
  \item Define occupancy mask: \(\texttt{occMask} = (\texttt{nodeAt} > 0)\).
  \item Build \texttt{componentAt} by projecting dense connected-component IDs from
  \texttt{conncomp(fullGraph.graph)} onto grid cells.
  \item Build \texttt{boundaryMask} on the grid from dense \texttt{boundary\_mask}.
  \item Fallback: if no boundary seed was provided, detect geometric perimeter pixels
  (4-neighbor exposure to outside).
\end{enumerate}

\subsection*{6. Boundary distance, radius, and thickness maps}
\paragraph{Simple}
Compute how far each occupied cell is from the boundary; convert this to local radius and thickness.

\paragraph{Detailed}
\begin{enumerate}
  \item Compute \texttt{boundaryStepDistance} using multi-source BFS on occupied cells
  (8-neighborhood, unit step count).
  \item Compute physical radius map with weighted propagation:
  \begin{itemize}
    \item initialize boundary seeds with radius 0,
    \item repeatedly expand to 8-neighbors inside occupancy,
    \item edge cost is physical distance
    \(\sqrt{(\Delta x)^2 + (\Delta y)^2}\) between grid coordinates.
  \end{itemize}
  \item Set thickness map on occupied cells:
  \[
  t(r,c) = 2\,r(r,c).
  \]
  \item Attach per-node values back to full-graph nodes using each node's \((xBin,yBin)\).
\end{enumerate}

Also store spacing statistics \texttt{dx}, \texttt{dy}, and \texttt{h} (medians of positive axis spacings).

\subsection*{7. Skeletonization with Zhang-Suen and island enforcement}
\paragraph{Simple}
Each connected occupied component is thinned to a one-pixel centerline skeleton.

\paragraph{Detailed}
For each component mask:
\begin{enumerate}
  \item Apply Zhang-Suen thinning until convergence (no pixel deletions in both sub-iterations).
  \item Pixel deletion conditions:
  \[
  2 \le N(p_1) \le 6, \quad S(p_1)=1,
  \]
  with standard two-pass logical constraints:
  \[
  \neg(p_2p_4p_6), \ \neg(p_4p_6p_8)
  \]
  in pass 1, and
  \[
  \neg(p_2p_4p_8), \ \neg(p_2p_6p_8)
  \]
  in pass 2.
  \item Enforce minimum skeleton nodes per occupancy component:
  if thinned skeleton count is below \texttt{MinIslandNodes}, force at least one representative pixel.
  \item Representative pixel selection:
  choose the component cell with maximum boundary-step distance
  (or deterministic fallback if all distances are infinite).
\end{enumerate}

\subsection*{8. Skeleton graph construction (\texttt{build\_skeleton\_graph\_from\_mask})}
\paragraph{Simple}
Convert skeleton pixels to graph nodes and connect neighboring pixels.

\paragraph{Detailed}
\begin{enumerate}
  \item Skeleton pixels become nodes; each stores:
  \begin{itemize}
    \item XY coordinates (from mapped full node),
    \item full node label and full node index,
    \item boundary step distance,
    \item node radius and node thickness.
  \end{itemize}
  \item Build edges by scanning four directions
  \((0,1),(1,-1),(1,0),(1,1)\) to avoid duplicate undirected insertion.
  \item Add edge only when:
  \begin{itemize}
    \item both pixels are skeleton pixels,
    \item both belong to the same connected component ID.
  \end{itemize}
  \item Additional diagonal rule:
  a diagonal link is allowed only if at least one orthogonal corner-support
  pixel belongs to the same component.
  \item Edge length is Euclidean distance between endpoint XY coordinates.
  \item Node role by skeleton degree:
  \begin{itemize}
    \item degree 0: \texttt{island}
    \item degree 1: \texttt{endpoint}
    \item degree 2: \texttt{chain}
    \item degree $>2$: \texttt{junction}
  \end{itemize}
\end{enumerate}

\subsection*{9. Skeleton to structural compression (\texttt{compress\_skeleton\_graph})}
\paragraph{Simple}
Trace paths through degree-2 chains and replace many small skeleton edges with fewer structural edges.

\paragraph{Detailed}
\begin{enumerate}
  \item Build edge lookup table and visited flags for skeleton edges.
  \item Process each skeleton connected component.
  \item Define key nodes in a component:
  \[
  \{v : \deg(v)\neq 2\}.
  \]
  \item If component has no key nodes, it is a pure cycle:
  trace a cycle path and process as closed path.
  \item Otherwise, for each key node and each neighbor:
  \begin{itemize}
    \item if edge not visited, trace chain with \texttt{trace\_chain},
    \item mark chain edges as visited,
    \item process as open or closed path.
  \end{itemize}
  \item After that, process any unvisited edges (safety pass).
\end{enumerate}

\paragraph{Open-path anchor rule}
For path length \(n\):
\[
m = n-2, \quad a = \min\big(m,\ \mathrm{round}(\texttt{DetailLevel}\cdot m)\big).
\]
Anchors are:
\begin{itemize}
  \item first node,
  \item \(a\) internal evenly-spaced positions,
  \item last node.
\end{itemize}

\paragraph{Closed-path anchor rule}
For cycle size \(n\):
\[
k=\min\left(n,\ \max\left(3,\ 1+\mathrm{round}(\texttt{DetailLevel}(n-1))\right)\right).
\]
At least 3 anchors are kept for every cycle.

\paragraph{Segment creation}
Each consecutive anchor pair creates one structural edge storing:
\begin{itemize}
  \item endpoint indices in reduced graph,
  \item \texttt{edge\_length} as polyline arc length,
  \item \texttt{edge\_polyline} as ordered skeleton coordinates,
  \item \texttt{edge\_full\_node\_paths} as ordered full-graph node indices,
  \item component ID.
\end{itemize}

\paragraph{Reduced node bookkeeping}
Reduced nodes are created lazily (\texttt{ensure\_reduced\_node}) and keep mappings:
\begin{itemize}
  \item structural node \(\rightarrow\) skeleton node index,
  \item structural node \(\rightarrow\) full node index.
\end{itemize}
Role promotion uses precedence:
\[
\texttt{waypoint} < \texttt{island} < \texttt{endpoint} < \texttt{junction}.
\]

\paragraph{Duplicate open-path guard}
If an open path would create a duplicate endpoint pair and has no internal anchor,
the code forces a middle anchor to split that path and avoid immediate duplicate edge insertion.

\subsection*{10. Structural cleanup (\texttt{cleanup\_structural\_graph})}
\paragraph{Simple}
After compression, run local merge rules to remove tiny artifacts.

\paragraph{Detailed}
\begin{enumerate}
  \item Estimate characteristic spacing \(h\):
  \begin{itemize}
    \item median positive axis spacing from node coordinates, or
    \item fallback median edge length.
  \end{itemize}
  \item Set cleanup thresholds:
  \[
  \ell_{\min}=1.5h,\quad r_{\text{merge}}=1.5h,\quad \text{maxIterations}=4.
  \]
  \item Initialize mutable cleanup state with:
  node fields, edge fields, path fields, and membership sets.
  \item Iterative pass (up to 4 rounds):
  \begin{enumerate}
    \item Repeated short-edge contractions:
    find eligible edge shorter than \(\ell_{\min}\), merge its endpoint pair.
    \item Repeated junction-cluster merges:
    find nearest same-component pair of degree \(>2\) nodes within \(r_{\text{merge}}\), merge.
  \end{enumerate}
  \item Stop early if an iteration makes no changes.
\end{enumerate}

\paragraph{Short-edge candidate eligibility}
For edge \((u,v)\), code skips candidates where:
\begin{itemize}
  \item \(u=v\),
  \item \(\deg(u)=1\) or \(\deg(v)=1\).
\end{itemize}
Candidate is accepted when:
\begin{itemize}
  \item \(\deg(u)\neq 2\) or \(\deg(v)\neq 2\), or
  \item \(u,v\) share a common neighbor.
\end{itemize}

\paragraph{Node merge mechanics}
When merging a node group:
\begin{enumerate}
  \item Choose representative node by:
  \begin{enumerate}
    \item highest degree,
    \item then highest thickness proxy,
    \item then smallest total XY distance to other candidates (medoid).
  \end{enumerate}
  \item Union membership metadata:
  merged full-node indices, merged skeleton-node indices, merged original structural indices.
  \item Rewire all incident edges from removed nodes to representative.
  \item Update edge polyline endpoints to representative coordinates.
  \item Compact state:
  remap node indices, remove deleted nodes, remove invalid edges, remove self loops.
  \item Deduplicate parallel edges by undirected pair and keep the longest.
\end{enumerate}

\paragraph{Cleanup outputs}
The final graph stores:
\begin{itemize}
  \item merged-membership arrays for traceability,
  \item updated node table columns (FullNodeIndex, SkeletonNodeIndex, Radius, Thickness, NodeRole, merged sets),
  \item \texttt{cleanup\_info} with thresholds and merge counters.
\end{itemize}

\subsection*{11. Structural edge thickness summaries}
\paragraph{Simple}
Each structural edge gets thickness statistics sampled along its stored full-node path.

\paragraph{Detailed}
For each structural edge path \(P_e\):
\[
t_{\min}(e)=\min_{i\in P_e} t_i,\quad
t_{\text{mean}}(e)=\frac{1}{|P_e|}\sum_{i\in P_e} t_i,\quad
t_{\max}(e)=\max_{i\in P_e} t_i.
\]
Then:
\[
t_{\text{bottleneck}}(e)=t_{\min}(e),\quad
\Delta t(e)=t_{\max}(e)-t_{\min}(e).
\]

These are written into:
\begin{itemize}
  \item \texttt{edge\_thickness\_min}
  \item \texttt{edge\_thickness\_mean}
  \item \texttt{edge\_thickness\_max}
  \item \texttt{edge\_thickness\_bottleneck}
  \item \texttt{edge\_thickness\_delta}
\end{itemize}

\subsection*{12. Persisted output structure (\texttt{sample.mat})}
\paragraph{Simple}
Each run produces one \texttt{sample.mat} with final graph and summary metadata.

\paragraph{Detailed}
\texttt{sample.mat} contains:
\begin{itemize}
  \item \texttt{sample.id}: run name
  \item \texttt{sample.source.inp\_path}: source INP path
  \item \texttt{sample.graph}: structural graph (after cleanup and thickness attachment)
  \item \texttt{sample.graph\_info}:
  \begin{itemize}
    \item mode and extraction parameters,
    \item node/edge counts for full, skeleton, and structural graphs.
  \end{itemize}
\end{itemize}

\section{Compact End-to-End Pseudocode}
\begin{verbatim}
for each run folder:
    inp = read_abaqus_inp(sheet_shell.inp)
    region = select_region_elements(inp, elset/pattern options)
    full = build_full_reference_graph(region)

    grid = rasterize(full.node_coords, full.boundary_mask, full.components)
    maps = compute_boundary_distance_and_thickness(grid)
    full = attach_full_graph_thickness(full, maps)

    skel_mask = thinning_per_component(grid.occMask, MinIslandNodes)
    skel = build_skeleton_graph_from_mask(skel_mask, full, maps)

    struct = compress_skeleton_graph(skel, DetailLevel)
    struct = cleanup_structural_graph(struct, skel)
    struct = attach_structural_edge_thickness(struct, full)

    save sample.mat with struct graph and graph_info
\end{verbatim}

\end{document}
